\section{Frobenius framing and a depth-one rapid-decay target}
\label{sec:frobenius-framing}

The additive normalization freedom visible in
\eqref{eq:tautological-recovery} and in the Kummer tower is often expressed
by the unipotent matrices
\[
 S_t=\begin{pmatrix}1&t\\0&1\end{pmatrix}.
\]
Inside an unframed Tannakian category, adding further tensors already fixed
by the motivic Galois group cannot remove this direction.  A Frobenius
operator with distinct graded eigenvalues does remove the corresponding
centralizer, but this fact has to be interpreted carefully: after forgetting
the filtration and comparison data, the same separation of slopes also makes
the local extension split.  We begin by recording both sides of this linear
algebra.

\begin{lemma}[Frobenius-framed rigidity]
\label{lem:frobenius-framed-rigidity}
Let $K$ be a field of characteristic zero and let $V$ be a two-dimensional
$K$-realization with a $K$-linearized Frobenius operator (for example, an
$f$th iterate of a semilinear Frobenius for which
$\sigma^f=\operatorname{id}_K$) which, in a basis compatible with its two
graded pieces, has matrix
\begin{equation}\label{eq:frobenius-matrix}
 \Phi_p=\begin{pmatrix}1&c_p\\0&q_p\end{pmatrix},
 \qquad q_p\ne1.
\end{equation}
Then the unipotent subgroup $U=\{S_t:t\in K\}$ has trivial intersection
with the centralizer of $\Phi_p$:
\begin{equation}\label{eq:frobenius-centralizer}
 U\cap Z_{\mathrm{GL}(V)}(\Phi_p)=\{1\}.
\end{equation}
In particular this applies when $q_p=p^f$ or $p^{-f}$, including $f=1$.
\end{lemma}

\begin{proof}
Direct multiplication gives
\[
 S_t\Phi_p-\Phi_pS_t
 =\begin{pmatrix}0&(q_p-1)t\\0&0\end{pmatrix}.
\]
Since $q_p\ne1$, commutation is equivalent to $t=0$.
\end{proof}

For a raw $\sigma$-semilinear Frobenius, the corresponding centralizer
condition is $q_pt=\sigma(t)$ rather than $(q_p-1)t=0$.  The associated
inhomogeneous equation has the opposite-looking consequence.

\begin{theorem}[Automatic splitting of a bare Frobenius extension]
\label{thm:semilinear-frobenius-splitting}
Let $K$ be a complete non-archimedean field, let
$\sigma\colon K\to K$ be an isometric automorphism, and let
$q\in K^\sigma$ satisfy $|q|\ne1$.  For every $c\in K$, consider the
$\sigma$-semilinear operator whose matrix is
\[
 A=\begin{pmatrix}1&c\\0&q\end{pmatrix}.
\]
There is a unique $t\in K$ such that
\begin{equation}\label{eq:semilinear-splitting-equation}
 qt-\sigma(t)=c.
\end{equation}
For this value of $t$ one has
\begin{equation}\label{eq:semilinear-diagonalization}
 S_t^{-1}A\,\sigma(S_t)
 =\begin{pmatrix}1&0\\0&q\end{pmatrix}.
\end{equation}
Thus the underlying unfiltered semilinear Frobenius extension is split, and
its upper-unipotent splitting is unique.
\end{theorem}

\begin{proof}
If $|q|<1$, the convergent series
\[
 t=-\sum_{n\ge0}q^n\sigma^{-n-1}(c)
\]
satisfies \eqref{eq:semilinear-splitting-equation}.  If $|q|>1$, the
corresponding solution is
\[
 t=\sum_{n\ge0}q^{-n-1}\sigma^n(c).
\]
Convergence follows because $\sigma$ is isometric.  If $t_1$ and $t_2$ are
two solutions, their difference $u$ satisfies $qu=\sigma(u)$; taking
absolute values gives $|q||u|=|u|$, and hence $u=0$.  Finally,
\[
 S_t^{-1}A\,\sigma(S_t)
 =\begin{pmatrix}1&c+\sigma(t)-qt\\0&q\end{pmatrix},
\]
which proves \eqref{eq:semilinear-diagonalization}.
\end{proof}

Theorem~\ref{thm:semilinear-frobenius-splitting} does not contradict
Lemma~\ref{lem:frobenius-framed-rigidity}.  The homogeneous centralizer
equation has only the zero solution, while the inhomogeneous splitting
equation has exactly one solution.  Consequently a marked Frobenius basis
may remove a gauge ambiguity, but the bare Frobenius isocrystal does not by
itself retain the global extension class.

\begin{proposition}[A Kummer local-to-global counterexample]
\label{prop:kummer-bare-frobenius-counterexample}
Let
\[
 M_2=[\mathbb Z\xrightarrow{1\mapsto2}\mathbb G_m]
\]
be the Kummer $1$-motive over $\mathbb Q$.  Its canonical sequence
\begin{equation}\label{eq:kummer-global-extension}
 0\longrightarrow[0\to\mathbb G_m]
 \longrightarrow M_2
 \longrightarrow[\mathbb Z\to0]
 \longrightarrow0
\end{equation}
is non-split in the isogeny category.  Nevertheless, $M_2$ has good
reduction at every prime $p\ne2$, and at each such prime its underlying
rational crystalline $F$-isocrystal is split.  Hence the implication
\[
 \left.
 \begin{array}{c}
 \text{the unfiltered rational $F$-isocrystal is split}\\
 \text{at every good prime}
 \end{array}
 \right\}
 \quad\Longrightarrow\quad
 \text{the global motive is split}
\]
is false.
\end{proposition}

\begin{proof}
A splitting of \eqref{eq:kummer-global-extension} up to isogeny would make
the point $2\in\mathbb G_m(\mathbb Q)$ torsion, equivalently would force
$2^m=1$ for some $m\ge1$.  Thus the global extension is non-split.  For
$p\ne2$, the point $2$ extends to a unit over $\mathbb Z_p$, so the motive
has good reduction.  Its reduction $\overline{2}\in\mathbb F_p^\times$ has
finite order.  The reduced Kummer $1$-motive therefore splits up to
isogeny, and applying the rational crystalline realization gives a split
sequence with graded slopes $0$ and $1$, up to the convention for signs.
This is the root-of-unity case of the Kummer calculation in
\cite[Example~6.3]{MohajerSebbar2025}; the surrounding slope formalism may be
compared with \cite{KedlayaIsocrystals}.
\end{proof}

The qualifier \emph{unfiltered} in the proposition is essential.  Neither
the $p$-adic Galois realization nor the filtered Frobenius module is asserted
to split.  The relative position of the slope splitting and the Hodge--de
Rham filtration can retain the Kummer extension; in the preceding example
its coordinate is represented, subject to conventions, by $\log_p2$.
Accordingly, a viable arithmetic enhancement of the Euler extension must
use filtered-Frobenius, $p$-adic integration, or syntomic comparison data,
not the bare $F$-isocrystal alone.

The rapid-decay pairing itself is well established: Bloch and Esnault
construct it for irregular connections on curves, while Hien develops the
surface and arbitrary-dimensional theories
\cite{BlochEsnault2004,Hien2007,Hien2009}.  The classical arithmetic
comparison is also instructive.  Huber and W\"ustholz determine the
algebraic linear relations among periods of classical $1$-motives
\cite{HuberWustholz}.  For $1$-motives over number fields with good
reduction, Mohajer and Sebbar construct canonical rational structures and a
$p$-adic analytic-subgroup framework at depths one and two
\cite{MohajerSebbar2025}.  On the irregular side, Snodgrass constructs
rational rapid-decay realizations and $E$-periods in the $E$-operator
setting under his stated hypotheses \cite{Snodgrass2026}.  None of these
results supplies the numerical-to-motivic implication needed below.

Let $M_\gamma$ denote the Euler--Mascheroni exponential motive in the
normalization of Fres\'an and Jossen.  They prove that it is a non-split
extension
\begin{equation}\label{eq:euler-motive-extension}
 0\longrightarrow\mathbb Q(0)
 \longrightarrow M_\gamma
 \longrightarrow\mathbb Q(-1)
 \longrightarrow0,
\end{equation}
that $\operatorname{Hom}(M_\gamma,\mathbb Q(0))=0$, and that its complex
extension coordinate is $\gamma$ \cite[Sec.~12.8]{FresanJossen}.

The relevant motivic separation from classical periods is already known;
it is not an additional mixed-extension conjecture.

\begin{theorem}[Fres\'an--Jossen classical separation]
\label{thm:fresan-jossen-classical-separation}
Let $N$ be any classical motive over $\mathbb Q$, put
$N^+=N\oplus M_\gamma$, and denote their motivic Galois groups by
$G_N$ and $G_{N^+}$.  Then the natural faithfully flat morphism
\[
 G_{N^+}\longrightarrow G_N
\]
has a subgroup isomorphic to $\mathbb G_a$ in its kernel.  In particular,
\begin{equation}\label{eq:fresan-jossen-dimension-jump}
 \dim G_{N^+}>\dim G_N.
\end{equation}
The statement applies, in particular, to every finite direct sum of Tate
and Kummer motives.
\end{theorem}

\begin{proof}
Fres\'an and Jossen show that the perverse realization of a classical
motive has trivial fundamental group, whereas adjoining $M_\gamma$
introduces the full unipotent group $\mathbb G_a$.  Their comparison diagram
embeds this copy of $\mathbb G_a$ into $G_{N^+}$ and maps it trivially to
$G_N$.  This is precisely the proof of
\cite[Cor.~12.8.12]{FresanJossen}; see also their
Proposition~12.8.7 and Corollary~12.8.8.
\end{proof}

Theorem~\ref{thm:fresan-jossen-classical-separation} settles the motivic
part of the Euler--Kummer separation problem, object by object for every
finite classical family.  What remains open is the comparison step that
turns a numerical relation among complex period coordinates into a motivic
relation.  The minimal form needed for $\gamma$ is the following degree-one
fullness statement.

\begin{conjecture}[Euler rapid-decay comparison fullness]
\label{conj:rapid-decay-analytic-subgroup}
Assume that $M_\gamma$ is equipped with its rational rapid-decay realization
and with coherent filtered-Frobenius or syntomic comparison data at the good
places.  If
\begin{equation}\label{eq:euler-splitting-relation}
 \lambda\gamma+\beta=0,
 \qquad \lambda,\beta\in\Qbar,\quad \lambda\ne0,
\end{equation}
then this numerical relation is induced by a morphism in the
exponential-motivic category; equivalently in the present rank-two setting,
it forces a splitting of \eqref{eq:euler-motive-extension}.
\end{conjecture}

This is deliberately a rank-two, degree-one specialization of the
exponential-period philosophy.  It is not a consequence of the existence
of rapid-decay homology, a rational structure, or bare Frobenius
realizations.  The filtered or syntomic data are possible tools for proving
the comparison assertion, not substitutes for it.

\begin{proposition}[Conditional transcendence]
\label{prop:rapid-decay-implies-gamma}
Conjecture~\ref{conj:rapid-decay-analytic-subgroup} implies that $\gamma$ is
transcendental.
\end{proposition}

\begin{proof}
If $\gamma$ were algebraic, then
\eqref{eq:euler-splitting-relation} would hold with $\lambda=1$ and
$\beta=-\gamma$.  The conjecture would force
\eqref{eq:euler-motive-extension} to split, contradicting the non-splitting
proved by Fres\'an and Jossen
\cite[Prop.~12.8.7]{FresanJossen}.
\end{proof}

\begin{conjecture}[Euler--Baker linear period theorem]
\label{conj:euler-baker}
Let $\alpha_1,\ldots,\alpha_r\in\Qbar^\times$ be multiplicatively
independent, and fix branches of their logarithms.  If
\begin{equation}\label{eq:euler-baker-relation}
 \lambda\gamma+\beta+\nu(2\pi i)
 +\sum_{j=1}^r\mu_j\Log\alpha_j=0,
 \qquad
 \lambda,\beta,\nu,\mu_j\in\Qbar,
\end{equation}
then
\[
 \lambda=\beta=\nu=\mu_1=\cdots=\mu_r=0.
\]
\end{conjecture}

For clarity, let $N_{\boldsymbol\alpha}$ be a finite classical motive
containing the relevant Tate and Kummer objects, so that its period matrix
contains $1$, $2\pi i$, and the chosen $\Log\alpha_j$.

\begin{proposition}[The exact relative comparison gap]
\label{prop:relative-fullness-implies-euler-baker}
Assume degree-one comparison fullness for
$N_{\boldsymbol\alpha}\oplus M_\gamma$: every $\Qbar$-linear relation among
the displayed comparison-matrix entries is induced by motivic
functoriality.  Then Conjecture~\ref{conj:euler-baker} holds for
$\alpha_1,\ldots,\alpha_r$.
\end{proposition}

\begin{proof}
Under comparison fullness, a relation
\eqref{eq:euler-baker-relation} is motivic.  The copy of $\mathbb G_a$ in
the kernel furnished by
Theorem~\ref{thm:fresan-jossen-classical-separation} fixes all period
coordinates coming from $N_{\boldsymbol\alpha}$ and translates the Euler
extension coordinate.  Invariance of the relation therefore forces
$\lambda=0$.  The remaining relation
\[
 \beta+\nu(2\pi i)+\sum_{j=1}^r\mu_j\Log\alpha_j=0
\]
is classical.  The Baker--Lindemann theorem \cite{Baker1975}, together with the assumed
multiplicative independence, gives
$\beta=\nu=\mu_1=\cdots=\mu_r=0$.
\end{proof}

Thus no additional theorem asserting motivic independence of the Euler
extension from the span of Kummer extensions is missing: that separation is
already contained in
Theorem~\ref{thm:fresan-jossen-classical-separation}.  The genuine missing
input for the Euler--Baker statement is the relative
\emph{numerical-to-motivic comparison theorem} used in
Proposition~\ref{prop:relative-fullness-implies-euler-baker}.

\begin{remark}[An adelic formulation of the missing estimate]
Let $K$ now be a number field.  The local Dwork scales
$p^{-1/(p-1)}$ should not be identified numerically across different
primes.  A viable slope strategy would first have to place coherent
filtered-Frobenius or syntomic norms on one rapid-decay extension and then
form an adelic degree
\[
 \widehat{\deg}(Ks)=-\sum_{w\in M_K}n_w\log\|s\|_w.
\]
It would additionally require a theorem showing that an algebraic complex
period relation produces a global section of impossible adelic degree
unless the extension splits.  The bare slope decomposition cannot provide
this: Proposition~\ref{prop:kummer-bare-frobenius-counterexample} shows
that it can split at every good place while the global motive remains
non-split.  Neither the required coherence of the filtered local norms nor
the global slope inequality is presently proved.
\end{remark}

The conclusion is diagnostic.  Frobenius centralizer rigidity is a correct
uniqueness statement, but the rank-two scalar semilinear model of
Theorem~\ref{thm:semilinear-frobenius-splitting} splits automatically when
\(|q|\ne1\).  Fres\'an and Jossen already prove that the
Euler $\mathbb G_a$-direction survives after adjoining every fixed
classical motive.  The remaining depth-one problem is therefore exactly a
comparison/fullness theorem transferring a complex rapid-decay relation to
the motivic category.  Its minimal rank-two case would prove the
transcendence of $\gamma$; its finite relative Kummer form would prove the
Euler--Baker linear period theorem.
